% !TeX root = ../main.tex

% Библиографический список
\begin{ltrtr}

	\bibitem{ogawa1990}
	\textit{Ogawa~S., Lee~T.\,M., Kay~A.\,R., Tank~D.\,W.} Brain magnetic resonance imaging with contrast dependent on blood oxygenation~// Proc. Natl. Acad. Sci. USA. 1990. Vol.~87, no.~24. P.~9868--9872. \href{https://doi.org/10.1073/pnas.87.24.9868}{DOI:~10.1073/\allowbreak pnas.87.24.9868}.

	\bibitem{kwong1992}
	\textit{Kwong~K.\,K., Belliveau~J.\,W., Chesler~D.\,A., et al.} Dynamic magnetic resonance imaging of human brain activity during primary sensory stimulation~// Proc. Natl. Acad. Sci. USA. 1992. Vol.~89, no.~12. P.~5675--5679. \href{https://doi.org/10.1073/pnas.89.12.5675}{DOI:~10.1073/\allowbreak pnas.89.12.5675}.

	\bibitem{kireev2013}
	\textit{Киреев~М.\,В., Захс~Д.\,В., Коротков~А.\,Д., Медведев~С.\,В.} Современные методы функциональной томографической нейровизуализации в исследовании функций больного и здорового мозга~// Рос. физиол. журн. им. И.\,М. Сеченова. 2013. Т.~99, №~1. С.~53--62.

	\bibitem{myznikov2026}
	\textit{Мызников~А.\,Д., Селиверстов~Р.\,Ю., Коротков~А.\,Д., Алиев~Г.\,Б., Чередниченко~Д.\,В., Дидур~М.\,Д., Киреев~М.\,В.} Негативный BOLD-сигнал как способ комплиментарной характеристики мозгового обеспечения функций: потенциал использования в предхирургической подготовке методом функциональной магнитно-резонансной томографии~// Рос. физиол. журн. им. И.\,М. Сеченова. 2026. Т.~112, №~6. С.~1337--1373. \href{https://doi.org/10.7868/S2658655X26060014}{DOI:~10.7868/\allowbreak S2658655X26060014}.

	\bibitem{friston2000}
	\textit{Friston~K.\,J., Mechelli~A., Turner~R., Price~C.\,J.} Nonlinear responses in fMRI: the Balloon model, Volterra kernels, and other hemodynamics~// NeuroImage. 2000. Vol.~12, no.~4. P.~466--477. \href{https://doi.org/10.1006/nimg.2000.0630}{DOI:~10.1006/\allowbreak nimg.2000.0630}.

	\bibitem{buxton2004}
	\textit{Buxton~R.\,B., Uludag~K., Dubowitz~D.\,J., Liu~T.\,T.} Modeling the hemodynamic response to brain activation~// NeuroImage. 2004. Vol.~23, Suppl.~1. P.~S220--S233. \href{https://doi.org/10.1016/j.neuroimage.2004.07.013}{DOI:~10.1016/\allowbreak j.neuroimage.2004.07.013}.

	\bibitem{mandeville1999}
	\textit{Mandeville~J.\,B., Marota~J.\,J.\,A., Ayata~C., et al.} Evidence of a cerebrovascular postarteriole windkessel with delayed compliance~// J. Cereb. Blood Flow Metab. 1999. Vol.~19, no.~6. P.~679--689. \href{https://doi.org/10.1097/00004647-199906000-00012}{DOI:~10.1097/\allowbreak 00004647-\allowbreak 199906000-\allowbreak 00012}.

	\bibitem{masharipov2025}
	\textit{Машарипов~Р.\,С., Дидур~М.\,Д., Чередниченко~Д.\,В., Киреев~М.\,В.} Изучение контекстно-зависимых функциональных связей с помощью функциональной магнитно-резонансной томографии~// Журн. высш. нервн. деят. им. И.\,П. Павлова. 2025. Т.~75, №~5. С.~530--552. \href{https://doi.org/10.31857/S0044467725050037}{DOI:~10.31857/\allowbreak S0044467725050037}.

	\bibitem{friston1998b}
	\textit{Friston~K.\,J., Fletcher~P., Josephs~O., et al.} Event-related fMRI: characterizing differential responses~// NeuroImage. 1998. Vol.~7, no.~1. P.~30--40. \href{https://doi.org/10.1006/nimg.1997.0306}{DOI:~10.1006/\allowbreak nimg.1997.0306}.

	\bibitem{friston1994}
	\textit{Friston~K.\,J., Holmes~A.\,P., Worsley~K.\,J., et al.} Statistical parametric maps in functional imaging: a general linear approach~// Hum. Brain Mapp. 1994. Vol.~2, no.~4. P.~189--210. \href{https://doi.org/10.1002/hbm.460020402}{DOI:~10.1002/\allowbreak hbm.460020402}.

	\bibitem{medvedev2012}
	\textit{Медведев~С.\,В., Киреев~М.\,В., Коротков~А.\,Д.} Функциональная МРТ и ПЭТ в клинической практике~--- общие принципы и ограничения при картировании функций мозга~// Лучевая диагностика и терапия. 2012. №~3(3). С.~20--26.

	\bibitem{monti2011}
	\textit{Monti~M.\,M.} Statistical analysis of fMRI time-series: a critical review of the GLM approach~// Front. Hum. Neurosci. 2011. Vol.~5. P.~28. \href{https://doi.org/10.3389/fnhum.2011.00028}{DOI:~10.3389/\allowbreak fnhum.2011.00028}.

	\bibitem{handwerker2012}
	\textit{Handwerker~D.\,A., Gonzalez-Castillo~J., D'Esposito~M., Bandettini~P.\,A.} The continuing challenge of understanding and modeling hemodynamic variation in fMRI~// NeuroImage. 2012. Vol.~62, no.~2. P.~1017--1023. \href{https://doi.org/10.1016/j.neuroimage.2012.02.015}{DOI:~10.1016/\allowbreak j.neuroimage.2012.02.015}.

	\bibitem{lindquist2009}
	\textit{Lindquist~M.\,A., Meng Loh~J., Atlas~L.\,Y., Wager~T.\,D.} Modeling the hemodynamic response function in fMRI: efficiency, bias and mis-modeling~// NeuroImage. 2009. Vol.~45, no.~1, Suppl.~1. P.~S187--S198. \href{https://doi.org/10.1016/j.neuroimage.2008.10.065}{DOI:~10.1016/\allowbreak j.neuroimage.2008.10.065}.

	\bibitem{polimeni2021}
	\textit{Polimeni~J.\,R., Lewis~L.\,D.} Imaging faster neural dynamics with fast fMRI: a need for updated models of the hemodynamic response~// Prog. Neurobiol. 2021. Vol.~207. P.~102174. \href{https://doi.org/10.1016/j.pneurobio.2021.102174}{DOI:~10.1016/\allowbreak j.pneurobio.2021.102174}.

	\bibitem{woolrich2004}
	\textit{Woolrich~M.\,W., Behrens~T.\,E.\,J., Smith~S.\,M.} Constrained linear basis sets for HRF modelling using Variational Bayes~// NeuroImage. 2004. Vol.~21, no.~4. P.~1748--1761. \href{https://doi.org/10.1016/j.neuroimage.2003.12.024}{DOI:~10.1016/\allowbreak j.neuroimage.2003.12.024}.

	\bibitem{wager2003}
	\textit{Wager~T.\,D., Nichols~T.\,E.} Optimization of experimental design in fMRI: a general framework using a genetic algorithm~// NeuroImage. 2003. Vol.~18, no.~2. P.~293--309. \href{https://doi.org/10.1016/S1053-8119(02)00046-0}{DOI:~10.1016/\allowbreak S1053-\allowbreak 8119(02)00046-\allowbreak 0}.

	\bibitem{friston1998}
	\textit{Friston~K.\,J., Josephs~O., Rees~G., Turner~R.} Nonlinear event-related responses in fMRI~// Magn. Reson. Med. 1998. Vol.~39, no.~1. P.~41--52. \href{https://doi.org/10.1002/mrm.1910390109}{DOI:~10.1002/\allowbreak mrm.1910390109}.

	\bibitem{robinson2006}
	\textit{Robinson~P.\,A., Drysdale~P.\,M., Van der Merwe~H., et al.} BOLD responses to stimuli: dependence on frequency, stimulus form, amplitude, and repetition rate~// NeuroImage. 2006. Vol.~31, no.~2. P.~585--599. \href{https://doi.org/10.1016/j.neuroimage.2005.12.026}{DOI:~10.1016/\allowbreak j.neuroimage.2005.12.026}.

	\bibitem{birn2001}
	\textit{Birn~R.\,M., Saad~Z.\,S., Bandettini~P.\,A.} Spatial heterogeneity of the nonlinear hemodynamic response~// NeuroImage. 2001. Vol.~14, no.~4. P.~817--826. \href{https://doi.org/10.1006/nimg.2001.0873}{DOI:~10.1006/\allowbreak nimg.2001.0873}.

	\bibitem{mechelli2001}
	\textit{Mechelli~A., Price~C.\,J., Friston~K.\,J.} Nonlinear coupling between evoked rCBF and BOLD signals: a simulation study of hemodynamic responses~// NeuroImage. 2001. Vol.~14, no.~4. P.~862--872. \href{https://doi.org/10.1006/nimg.2001.0876}{DOI:~10.1006/\allowbreak nimg.2001.0876}.

	\bibitem{magri2011}
	\textit{Magri~C., Logothetis~N.\,K., Panzeri~S.} Investigating static nonlinearities in neurovascular coupling~// Magn. Reson. Imaging. 2011. Vol.~29, no.~10. P.~1358--1364. \href{https://doi.org/10.1016/j.mri.2011.04.017}{DOI:~10.1016/\allowbreak j.mri.2011.04.017}.

	\bibitem{buxton1998}
	\textit{Buxton~R.\,B., Wong~E.\,C., Frank~L.\,R.} Dynamics of blood flow and oxygenation changes during brain activation: the Balloon model~// Magn. Reson. Med. 1998. Vol.~39, no.~6. P.~855--864. \href{https://doi.org/10.1002/mrm.1910390602}{DOI:~10.1002/\allowbreak mrm.1910390602}.

	\bibitem{friston2003}
	\textit{Friston~K.\,J., Harrison~L., Penny~W.} Dynamic causal modelling~// NeuroImage. 2003. Vol.~19, no.~4. P.~1273--1302. \href{https://doi.org/10.1016/S1053-8119(03)00202-7}{DOI:~10.1016/\allowbreak S1053-\allowbreak 8119(03)00202-\allowbreak 7}.

	\bibitem{wager2005}
	\textit{Wager~T.\,D., Vazquez~A., Hernandez~L., Noll~D.\,C.} Accounting for nonlinear BOLD effects in fMRI: parameter estimates and a model for prediction in rapid event-related studies~// NeuroImage. 2005. Vol.~25, no.~1. P.~206--218. \href{https://doi.org/10.1016/j.neuroimage.2004.11.008}{DOI:~10.1016/\allowbreak j.neuroimage.2004.11.008}.

	\bibitem{welvaert2014}
	\textit{Welvaert~M., Rosseel~Y.} A review of fMRI simulation studies~// PLoS ONE. 2014. Vol.~9, no.~7. P.~e101953. \href{https://doi.org/10.1371/journal.pone.0101953}{DOI:~10.1371/\allowbreak journal.pone.0101953}.

	\bibitem{ellis2020}
	\textit{Ellis~C.\,T., Baldassano~C., Schapiro~A.\,C., Cai~M.\,B., Cohen~J.\,D.} Facilitating open-science with realistic fMRI simulation: validation and application~// PeerJ. 2020. Vol.~8. P.~e8564. \href{https://doi.org/10.7717/peerj.8564}{DOI:~10.7717/\allowbreak peerj.8564}.

	\bibitem{cakan2023}
	\textit{Cakan~C., Jajcay~N., Obermayer~K.} neurolib: a simulation framework for whole-brain neural mass modeling~// Cogn. Comput. 2023. Vol.~15, no.~4. P.~1132--1152. \href{https://doi.org/10.1007/s12559-021-09931-9}{DOI:~10.1007/\allowbreak s12559-\allowbreak 021-\allowbreak 09931-\allowbreak 9}.

	\bibitem{havlicek2015}
	\textit{Havlicek~M., Roebroeck~A., Friston~K., Gardumi~A., Ivanov~D., Uludag~K.} Physiologically informed dynamic causal modeling of fMRI data~// NeuroImage. 2015. Vol.~122. P.~355--372. \href{https://doi.org/10.1016/j.neuroimage.2015.07.078}{DOI:~10.1016/\allowbreak j.neuroimage.2015.07.078}.

	\bibitem{sanzleon2015}
	\textit{Sanz-Leon~P., Knock~S.\,A., Spiegler~A., Jirsa~V.\,K.} Mathematical framework for large-scale brain network modeling in The Virtual Brain~// NeuroImage. 2015. Vol.~111. P.~385--430. \href{https://doi.org/10.1016/j.neuroimage.2015.01.002}{DOI:~10.1016/\allowbreak j.neuroimage.2015.01.002}.

	\bibitem{deneux2006}
	\textit{Deneux~T., Faugeras~O.} Using nonlinear models in fMRI data analysis: model selection and activation detection~// NeuroImage. 2006. Vol.~32, no.~4. P.~1669--1689. \href{https://doi.org/10.1016/j.neuroimage.2006.03.006}{DOI:~10.1016/\allowbreak j.neuroimage.2006.03.006}.

	\bibitem{volterra1930}
	\textit{Volterra~V.} Theory of Functionals and of Integral and Integro-Differential Equations. London: Blackie and Son, 1930.

	\bibitem{schetzen1980}
	\textit{Schetzen~M.} The Volterra and Wiener Theories of Nonlinear Systems. New York: Wiley, 1980.

	\bibitem{marmarelis2004}
	\textit{Marmarelis~V.\,Z.} Nonlinear Dynamic Modeling of Physiological Systems. Hoboken: IEEE Press; Wiley-Interscience, 2004. \href{https://doi.org/10.1002/9780471679370}{DOI:~10.1002/\allowbreak 9780471679370}.

	\bibitem{hunter1986}
	\textit{Hunter~I.\,W., Korenberg~M.\,J.} The identification of nonlinear biological systems: Wiener and Hammerstein cascade models~// Biol. Cybern. 1986. Vol.~55, no.~2--3. P.~135--144. \href{https://doi.org/10.1007/BF00341929}{DOI:~10.1007/\allowbreak BF00341929}.

	\bibitem{boyd1985}
	\textit{Boyd~S., Chua~L.\,O.} Fading memory and the problem of approximating nonlinear operators with Volterra series~// IEEE Trans. Circuits Syst. 1985. Vol.~32, no.~11. P.~1150--1161. \href{https://doi.org/10.1109/TCS.1985.1085649}{DOI:~10.1109/\allowbreak TCS.1985.1085649}.

	\bibitem{pintelon2012}
	\textit{Pintelon~R., Schoukens~J.} System Identification: A Frequency Domain Approach. 2nd ed. Hoboken: Wiley, 2012. \href{https://doi.org/10.1002/9781118287422}{DOI:~10.1002/\allowbreak 9781118287422}.

	\bibitem{schoukens2015}
	\textit{Schoukens~J., Pintelon~R., Rolain~Y., et al.} Structure discrimination in block-oriented models using linear approximations: a theoretic framework~// Automatica. 2015. Vol.~53. P.~225--234. \href{https://doi.org/10.1016/j.automatica.2014.12.045}{DOI:~10.1016/\allowbreak j.automatica.2014.12.045}.

	\bibitem{glover1999}
	\textit{Glover~G.\,H.} Deconvolution of impulse response in event-related BOLD fMRI~// NeuroImage. 1999. Vol.~9, no.~4. P.~416--429. \href{https://doi.org/10.1006/nimg.1998.0419}{DOI:~10.1006/\allowbreak nimg.1998.0419}.

\end{ltrtr}
